% ============================================================
%  6. fejezet, A megvalósított rendszer
% ============================================================
\chapter{A megvalósított rendszer}
\label{ch:implementacio}

A fejezet a KCMamba modell köré épített szoftvert tárgyalja:
a kereskedési motort, az irányítópultot, a visszatesztelő infrastruktúrát
és a Binance-adaptert. A rendszer hexagonális architektúrában (\emph{Ports \& Adapters})
készült; az elveket a~\ref{sec:system_arch}. alfejezet ismertette,
itt az egyes adapterek konkrét megvalósítására fókuszálok.

% ─────────────────────────────────────────────────────────────────────────────
\section{Megvalósítási technológiák}
\label{sec:impl_tech}
% ─────────────────────────────────────────────────────────────────────────────

A rendszer teljes egészében \textbf{Python 3.11}-ben íródott.
A főbb felhasznált könyvtárak és technológiák:

\begin{table}[H]
  \centering
  \small
  \begin{tabular}{|l|l|p{7cm}|}
    \hline
    \textbf{Réteg} & \textbf{Technológia} & \textbf{Felhasználás} \\
    \hline\hline
    Modell         & PyTorch~2.x          & KCMamba forward-pass, gradiensszámítás \\
    Adatfolyam     & asyncio / aiohttp    & Eseményvezérelt, nem-blokkoló I/O \\
    Tőzsde-adapter & python-binance       & Élő piaci adat, megbízás-végrehajtás \\
    Webes felület  & aiohttp (HTTP szerver) & Statikus HTML/JS, REST JSON végpontok \\
    Frontend       & Vanilla JS + Chart.js & Interaktív grafikonok, nincs build-lépés \\
    Tesztelés      & pytest + unittest.mock & Unit- és integrációs tesztek (backtester) \\
    Konfigurálás   & pydantic v2          & Szigorú típusú beállítások validálása \\
    Naplózás       & Python logging       & Rotating file + konzol handler \\
    \hline
  \end{tabular}
  \caption{A rendszer főbb technológiai komponensei és azok szerepe.}
  \label{tab:tech_stack}
\end{table}

A frontend szándékosan nem használ React-et vagy Vue-t:
minden oldal önálló, statikus HTML-fájl, amelyet az \texttt{aiohttp}
szerver közvetlenül kiszolgál. Nincs \texttt{npm} build-lépés,
Chart.js-sel mégis interaktív vizualizáció valósul meg.

% ─────────────────────────────────────────────────────────────────────────────
\section{Komponens-terv}
\label{sec:component_plan}
% ─────────────────────────────────────────────────────────────────────────────

A rendszer egyetlen Python-folyamatban fut; a komponensek
az alábbiak szerint kapcsolódnak egymáshoz:

\noindent\textbf{Fő folyamat} (\texttt{main.py})\\
Elindítja az \texttt{asyncio} eseményhurkot és inicializálja
a \texttt{TradingEngine}-t, az \texttt{MLService}-t és az adaptert.

\noindent\textbf{TradingEngine} (\texttt{core/application/engine.py})\\
Az üzleti logika középpontja. Az \texttt{InMemoryEventBus}-on keresztül
fogad \texttt{PriceEvent}-eket, sorban futtatja a pipeline-lépéseket
(Prediction $\to$ Signal $\to$ Risk $\to$ Execution), majd a
\texttt{IBrokerGatewayPort}-on adja ki a megbízást.

\noindent\textbf{aiohttp webes szerver} (\texttt{adapters/web/dashboard.py})\\
Ugyanabban az \texttt{asyncio} event-loopban fut a kereskedési hurokkal.
1~Hz-es \texttt{asyncio.Task}-ként WebSocket-streamet küld a böngészőknek;
REST végpontokon (\texttt{/api/signals}, \texttt{/api/equity} stb.)
keresztül is lekérdezhető az állapot.

\noindent\textbf{Binance WebSocket / REST} (külső)\\
A \texttt{BinanceFeed} aszinkron WebSocket-kapcsolaton fogad
OHLCV-adatot, a \texttt{BinanceBroker} REST API-n küld megbízásokat.
Hálózati megszakadás esetén automatikus újracsatlakozás.

\noindent\textbf{Fájlrendszer}\\
Modellsúlyok (\texttt{*.pt}), naplófájlok, CSV-exportok és a
\texttt{.env} konfigurációs fájl. Adatbázist a rendszer nem használ;
minden futásidei állapot in-memory adatstruktúrákban él.

\begin{table}[H]
  \centering
  \small
  \begin{tabular}{|l|l|l|l|}
    \hline
    \textbf{Komponens} & \textbf{Protokoll} & \textbf{Irány} & \textbf{Szerepkör} \\
    \hline\hline
    TradingEngine $\leftrightarrow$ EventBus   & in-process call & kétirányú & esemény-pub/sub \\
    TradingEngine $\to$ BinanceBroker          & REST HTTPS      & kimenő    & megbízás-küldés \\
    BinanceFeed $\leftarrow$ Binance           & WebSocket WSS   & bejövő    & árfolyam-stream \\
    DashboardServer $\leftrightarrow$ böngésző & HTTP + WS       & kétirányú & vizualizáció \\
    MLService $\leftrightarrow$ fájlrendszer   & fájl I/O        & kétirányú & súlyok be/ki \\
    \hline
  \end{tabular}
  \caption{A rendszer fizikai komponenseinek kapcsolatrendszere futásidőben.}
  \label{tab:component_plan}
\end{table}

% ─────────────────────────────────────────────────────────────────────────────
\section{A webalapú irányítópult}
\label{sec:dashboard}
% ─────────────────────────────────────────────────────────────────────────────

Az irányítópult hat nézetből áll; ezek a bal oldali navigációs sávon
keresztül érhetők el. A backend 1~Hz-es WebSocket-streammel frissíti
az egyes paneleket.

% · · · · · · · · · · · · · · · · · · · · · · · · · · · · · · · · · · · · ·
\subsection{Főoldal, Áttekintő nézet}
\label{subsec:dash_overview}
% · · · · · · · · · · · · · · · · · · · · · · · · · · · · · · · · · · · · ·

A főoldal (\texttt{index.html}) a legfontosabb információkat sűríti össze:
tőke-görbe (equity curve), nyitott pozíciók, legutóbbi kereskedési jelzések
és a modell aktuális Kalman-gain értéke.

\begin{figure}[H]
  \centering
  % ── KÉPERNYŐKÉP HELYE ──────────────────────────────────────────────────
  \fbox{\begin{minipage}[c][6cm][c]{0.92\textwidth}
    \centering
    \textit{[Képernyőkép: Főoldal, Áttekintő nézet]}\\[4pt]
    \small Ide kerül a \texttt{dashboard/index.html} oldal képernyőképe,\\
    amely a tőke-görbét, a nyitott pozíciókat és a Kalman-gain értéket mutatja.
  \end{minipage}}
  % ── KÉPERNYŐKÉP VÉGE ───────────────────────────────────────────────────
  \caption{Az irányítópult főoldala. A felső sávon a tőke-görbe, alatta a
           nyitott pozíciók és a legutóbbi jelzések láthatók.}
  \label{fig:dash_overview}
\end{figure}

% · · · · · · · · · · · · · · · · · · · · · · · · · · · · · · · · · · · · ·
\subsection{Kereskedési jelzések}
\label{subsec:dash_signals}
% · · · · · · · · · · · · · · · · · · · · · · · · · · · · · · · · · · · · ·

A jelzésnapló (\texttt{signals.html}) az összes generált \textsc{Buy}/\textsc{Sell}/\textsc{Hold}
jelzést jeleníti meg időbélyeggel, eszköznévvel, jelzés-erősséggel és a modell
aktuális előrejelzési értékével együtt.  Az oldal szűrhető szimbólum és
jelzéstípus szerint.

\begin{figure}[H]
  \centering
  \fbox{\begin{minipage}[c][6cm][c]{0.92\textwidth}
    \centering
    \textit{[Képernyőkép: Kereskedési jelzések, \texttt{signals.html}]}\\[4pt]
    \small A jelzésnapló táblázata időrendi sorrendben,
    szűrési és export-funkcióval.
  \end{minipage}}
  \caption{A kereskedési jelzések nézete. Minden sor egy generált jelzést
           ábrázol a jelzés erősségével és a döntéshozó modell állapotával.}
  \label{fig:dash_signals}
\end{figure}

% · · · · · · · · · · · · · · · · · · · · · · · · · · · · · · · · · · · · ·
\subsection{Modell diagnosztika}
\label{subsec:dash_model}
% · · · · · · · · · · · · · · · · · · · · · · · · · · · · · · · · · · · · ·

A modell diagnosztikai oldalon (\texttt{model.html}) az előrejelzési teljesítmény
valós idejű metrikái, gördülő MSE, MAE, hit-rate, láthatók, grafikonon
megjelenítve az előre jelzett és a tényleges értékeket.

\begin{figure}[H]
  \centering
  \fbox{\begin{minipage}[c][6cm][c]{0.92\textwidth}
    \centering
    \textit{[Képernyőkép: Modell diagnosztika, \texttt{model.html}]}\\[4pt]
    \small Előrejelzett vs.\ tényleges értékek, gördülő MSE-görbe,
    konfidencia-intervallum.
  \end{minipage}}
  \caption{A modell diagnosztikai nézete. A Chart.js-alapú panel valós időben
           frissül az előrejelzésekkel és eltérésekkel.}
  \label{fig:dash_model}
\end{figure}

% · · · · · · · · · · · · · · · · · · · · · · · · · · · · · · · · · · · · ·
\subsection{Kalman-szűrő diagnosztika}
\label{subsec:dash_kalman}
% · · · · · · · · · · · · · · · · · · · · · · · · · · · · · · · · · · · · ·

A Kalman-diagnosztikai nézet (\texttt{kla\_kalman.html}) a modell belső
állapotát mutatja: az adaptív $K$-gain, az $A$ felejtési faktor és a $R$
zajbecslő időbeli alakulása.  Ez a nézet hasznos annak ellenőrzésekor,
hogy a modell valóban reagál-e a piaci volatilitás változásaira.

\begin{figure}[H]
  \centering
  \fbox{\begin{minipage}[c][6cm][c]{0.92\textwidth}
    \centering
    \textit{[Képernyőkép: Kalman-szűrő diagnosztika, \texttt{kla\_kalman.html}]}\\[4pt]
    \small A $K$-gain, $A$-faktor és $R$-zajbecslő időbeli változása
    közepes és magas volatilitású piaci periódusban.
  \end{minipage}}
  \caption{A Kalman-szűrő diagnosztikai nézete. A $K$-gain csökkenése
           magas volatilitás esetén megerősíti az adaptív szűrési mechanizmus
           helyes működését.}
  \label{fig:dash_kalman}
\end{figure}

% · · · · · · · · · · · · · · · · · · · · · · · · · · · · · · · · · · · · ·
\subsection{Offline tanítás}
\label{subsec:dash_training}
% · · · · · · · · · · · · · · · · · · · · · · · · · · · · · · · · · · · · ·

A tanítási nézet (\texttt{training.html}) az offline betanítás menetét
monitorozza: epoch-veszteség-görbe, validációs MSE, becsült befejezési idő
és az aktuálisan futó konfiguráció paraméterei.

\begin{figure}[H]
  \centering
  \fbox{\begin{minipage}[c][6cm][c]{0.92\textwidth}
    \centering
    \textit{[Képernyőkép: Offline tanítás, \texttt{training.html}]}\\[4pt]
    \small Epoch-veszteség-görbe, validációs MSE élő frissítéssel,
    haladási sáv és konfiguráció-összefoglaló.
  \end{minipage}}
  \caption{Az offline tanítás irányítópultja. A tanítás teljes menete
           nyomon követhető, beleértve a korai leállás feltételét is.}
  \label{fig:dash_training}
\end{figure}

% · · · · · · · · · · · · · · · · · · · · · · · · · · · · · · · · · · · · ·
\subsection{Benchmark és modell-összehasonlítás}
\label{subsec:dash_benchmark}
% · · · · · · · · · · · · · · · · · · · · · · · · · · · · · · · · · · · · ·

A benchmark nézet (\texttt{ltsf\_benchmark.html}) az LSTM, Fair Mamba és
KCMamba eredményeit jeleníti meg interaktív hőtérképen és vonalas grafikon
formájában, különböző zajszintekre és stride-konfigurációkra lebontva.

\begin{figure}[H]
  \centering
  \fbox{\begin{minipage}[c][6cm][c]{0.92\textwidth}
    \centering
    \textit{[Képernyőkép: LSTM/Mamba összehasonlítás, \texttt{ltsf\_benchmark.html}]}\\[4pt]
    \small Interaktív hőtérkép és vonalgrafikonok az MSE-eredményekről
    különböző zajszinteken és adatkészleteken.
  \end{minipage}}
  \caption{A benchmark összehasonlítási nézet. Az interaktív szűrők lehetővé
           teszik az adatkészlet, zajszint és stride dinamikus változtatását.}
  \label{fig:dash_benchmark}
\end{figure}

% ─────────────────────────────────────────────────────────────────────────────
\section{A visszatesztelő keretrendszer}
\label{sec:backtest}
% ─────────────────────────────────────────────────────────────────────────────

A visszatesztelő (\texttt{adapters/testing/backtest.py}) az összes
kereskedési logikát ugyanazon \texttt{IBrokerGatewayPort} és
\texttt{IMarketFeedPort} interfészeken keresztül futtatja, mint az élő rendszer.
A visszateszteléshez szimulált eseménybusz (\texttt{InMemoryEventBus}) és
determinisztikus időforrás (\texttt{FixedTimeSource}) kerül felhasználásra,
így a tesztek reprodukálhatók.

A futtatás menete: a historikus OHLCV-adatot a rendszer betölti és normalizálja,
offline betanítja a modellt, majd percenkénti szimulációs loop-ban végighalad
a validációs perióduson: $0{,}1\%$ taker és $0{,}1\%$ maker díjjakkal.
Az eredményt tőkegörbe, Sharpe-mutató, maximális visszaesés és win-rate
formájában naplózza.

\begin{figure}[H]
  \centering
  \fbox{\begin{minipage}[c][5.5cm][c]{0.92\textwidth}
    \centering
    \textit{[Képernyőkép: Visszateszt eredmények, tőkegörbe és metrikák]}\\[4pt]
    \small Tőkegörbe (equity curve), Sharpe-mutató, maximális visszaesés
    és havi hozam-heatmap a visszatesztelési periódusra.
  \end{minipage}}
  \caption{Visszateszt eredmény-összesítő. A tőkegörbe mutatja a stratégia
           kumulatív teljesítményét az adott historikus perióduson.}
  \label{fig:backtest_results}
\end{figure}

% ─────────────────────────────────────────────────────────────────────────────
\section{Kockázatkezelés}
\label{sec:risk_management}
% ─────────────────────────────────────────────────────────────────────────────

A kockázatkezelési réteg (\texttt{core/domain/risk.py}) négy alapszabályt
kényszerít ki. Pozícióméret-korlát: egyetlen nyitás sem haladhatja meg a tőke
$10\%$-át. Aszimmetrikus stop-loss: hosszú pozíciónál $-2\%$, rövidnél
$-1{,}5\%$ belépési ártól. Ha egyszerre 5-nél több szimbólum lenne nyitott,
az újabb jelzés blokkolódik, és ha a napi veszteség meghaladja a $3\%$-ot,
a motor az összes pozíciót zárja.

Ezentúl csak akkor adódik ki megbízás, ha a KCMamba $K$-gain értéke
meghaladja a $K_\text{min} = 0{,}15$ küszöböt, ami azt jelzi, hogy a modell
elegendő bizalommal rendel előrejelzéshez az aktuális adatminőség alapján.

\begin{figure}[H]
  \centering
  \fbox{\begin{minipage}[c][5cm][c]{0.92\textwidth}
    \centering
    \textit{[Képernyőkép: Kockázati napló és pozícióméret-panel]}\\[4pt]
    \small A kockázatkezelő modul naplója: megkísérelt és elutasított
    jelzések, stop-loss eseménysorozat, napi veszteség-progress sáv.
  \end{minipage}}
  \caption{A kockázatkezelési modul áttekintő nézete. Piros sor jelzi
           az elutasított jelzéseket; zöld a végrehajtottakat.}
  \label{fig:risk_panel}
\end{figure}

% ─────────────────────────────────────────────────────────────────────────────
\section{Élő kereskedés Binance-szal}
\label{sec:live_trading}
% ─────────────────────────────────────────────────────────────────────────────

Az éles kereskedési adapter (\texttt{adapters/infrastructure/binance\_broker.py})
a Binance Futures REST és WebSocket végpontjait használja. Az árfolyam-frissítések
az \texttt{IEventBusPort}-on keresztül jutnak el a \texttt{TradingEngine}-hez;
az üzleti logika nem tud a Binance-specifikus adatformátumról; az adapter cserélhető.

\begin{figure}[H]
  \centering
  \fbox{\begin{minipage}[c][5cm][c]{0.92\textwidth}
    \centering
    \textit{[Képernyőkép: Élő kereskedés, Binance pozíciók és megbízások]}\\[4pt]
    \small Nyitott Binance Futures pozíciók, függőben lévő megbízások,
    tőke-frissítések valós időben.
  \end{minipage}}
  \caption{Az élő Binance-kereskedési nézet. A panel szinkronban van a
           Binance számlával; az egyenleg és a pozíciók másodpercenként
           frissülnek.}
  \label{fig:live_trading}
\end{figure}

A rendszer papírkereskedési (\emph{paper trading}) módban is futtatható:
ebben az esetben a Binance-adapter le van cserélve a
\texttt{adapters/testing/broker.py} szimulált végrehajtójára, így az algoritmus
valódi tőke kockáztatása nélkül tesztelhető éles piaci adaton.

% ─────────────────────────────────────────────────────────────────────────────
\section{A rendszer futtatása}
\label{sec:running}
% ─────────────────────────────────────────────────────────────────────────────

A teljes rendszer egyetlen belépési ponton, a \texttt{main.py} fájlon
keresztül indítható:

\begin{verbatim}
# Visszatesztelési mód (nincs hálózati kapcsolat)
python main.py --mode backtest --symbol BTCUSDT --from 2024-01-01 --to 2024-12-31

# Papírkereskedési mód (Binance feed, szimulált végrehajtás)
python main.py --mode paper --symbol BTCUSDT ETHUSDT

# Éles kereskedési mód (Binance API-kulcs szükséges)
python main.py --mode live --symbol BTCUSDT
\end{verbatim}

Az irányítópult automatikusan elindul, és a \texttt{http://localhost:8080}
címen érhető el.  Az API-kulcsok \texttt{.env} fájlból töltődnek be, sohasem
kerülnek a forráskódba.

\begin{figure}[H]
  \centering
  \fbox{\begin{minipage}[c][4.5cm][c]{0.92\textwidth}
    \centering
    \textit{[Képernyőkép: Terminál, rendszerindítás és naplósorok]}\\[4pt]
    \small A konzolkimenet a rendszerindítás során: modelltöltés, WebSocket
    csatlakozás megerősítése, eseményhurok indulása.
  \end{minipage}}
  \caption{A rendszer indítási naplója. Az egyes adapterek sorban inicializálódnak;
           hiba esetén a rendszer részletes hibaüzenetet és stack trace-t ír ki.}
  \label{fig:startup_log}
\end{figure}

% ─────────────────────────────────────────────────────────────────────────────
\section{Felhasználói dokumentáció}
\label{sec:user_docs}
% ─────────────────────────────────────────────────────────────────────────────

\subsection{A szoftver célja és célközönsége}

A rendszer kriptovaluta idősorok automatizált elemzésére és kereskedési
jelzések generálására készült, elsősorban BTCUSDT és ETHUSDT párokon.
Célközönsége fejlesztők és kutatók, akik ML-alapú kereskedési stratégiákat
szeretnének visszatesztelni vagy élőben futtatni.

\subsection{Rendszerkövetelmények}

\begin{table}[H]
  \centering
  \small
  \begin{tabular}{|l|l|l|}
    \hline
    \textbf{Összetevő} & \textbf{Minimum} & \textbf{Ajánlott} \\
    \hline\hline
    Operációs rendszer & Linux/macOS/Windows 10+ & Ubuntu 22.04 LTS \\
    Python verziója    & 3.10                   & 3.11 \\
    RAM                & 4~GB                   & 8~GB \\
    Lemezterület       & 2~GB                   & 5~GB \\
    GPU (opcionális)   & n.a.                      & NVIDIA $\ge$6~GB~VRAM \\
    Hálózat            & n.a. (backtest módban)    & Stabil internet (live) \\
    \hline
  \end{tabular}
  \caption{A rendszer minimális és ajánlott hardver/szoftvere.}
  \label{tab:sysreq}
\end{table}

\subsection{Telepítési útmutató}

\begin{enumerate}
  \item Klónozás és virtuális környezet létrehozása:
\begin{verbatim}
git clone <repo_url> && cd diplomamunkakod
python -m venv .venv && source .venv/bin/activate
\end{verbatim}
  \item Függőségek telepítése (ML-csomagok opcionálisak):
\begin{verbatim}
pip install -r requirements.txt
pip install -r requirements-ml.txt   # GPU-ra: CUDA-kompatibilis PyTorch
\end{verbatim}
  \item API-kulcsok beállítása (csak élő/paper módban szükséges):
\begin{verbatim}
cp .env.example .env
# Szerkeszd a .env fájlt: BINANCE_API_KEY, BINANCE_SECRET
\end{verbatim}
\end{enumerate}

\subsection{Használati útmutató}

\subsubsection*{Első indítás (paper-trading demó)}

A rendszer alap-használati módja egy előtérben futó paper-trading demó. A
\texttt{main.py} szkript végrehajtása egyszerre három háttérszolgáltatást
indít: egy offline tanítófolyamot, a Binance WebSocket-feed-et és a webes
irányítópultot.

\begin{verbatim}
source .venv/bin/activate
python main.py
\end{verbatim}

Sikeres indítás esetén a konzol az alábbi blokkokat írja ki: a \texttt{core}
modulok betöltése, az adapterek inicializálása, a historikus adatcsomag
letöltése (a \texttt{warmup\_start\_str} szerint, alapértelmezetten 3 napra
visszamenőleg), majd a
\texttt{[BOOT] TradingEngine started} szalag, amely után az irányítópult
a \texttt{http://127.0.0.1:8080} címen érhető el.

\subsubsection*{Dashboard-nézetek}

Az irányítópult felső navigációs sávja nyolc nézetre tagolódik, három
kategóriába rendezve:

\begin{table}[H]
  \centering
  \small
  \begin{tabular}{|l|l|p{7.2cm}|}
    \hline
    \textbf{Kategória} & \textbf{Oldal} & \textbf{Tartalom} \\
    \hline\hline
    Pénzügyi  & \texttt{/}              & Tőkegörbe, nyitott pozíciók, élő ár-chart \\
    Pénzügyi  & \texttt{/signals.html}  & Generált jelzések listája, kockázati státusszal \\
    Pénzügyi  & \texttt{/model.html}    & Modellverzió, kalibráció, diagnosztikai metrikák \\
    Pénzügyi  & \texttt{/training.html} & Offline tanítás előrehaladása, loss-görbe \\
    Benchmark & \texttt{/benchmark.html}     & Szintetikus zajos teszt eredménye \\
    Benchmark & \texttt{/ltsf\_benchmark.html}& LTSF (Exchange, ETTh) benchmark-eredmények \\
    KCA elemzés & \texttt{/kla\_kalman.html}   & Kalman-paraméterek ($K$, $A$, $R$) valós idejű görbéi \\
    KCA elemzés & \texttt{/complexity.html}    & Paraméterszám $\times$ teljesítmény heatmap \\
    \hline
  \end{tabular}
  \caption{Az irányítópult nézeteinek szerepköre.}
  \label{tab:dashboard_pages}
\end{table}

A jobb alsó sarokban elhelyezett \texttt{?} gombra kattintva oldal-specifikus
súgóablak nyílik, amely röviden elmagyarázza az adott nézetet. A lapok közti
átjárás a felső menüsávból történik; a böngésző \emph{vissza}-gombja is
használható.

\subsubsection*{Használati minták}

A rendszer négy tipikus forgatókönyvben futtatható. A paper-trading demó
(alapértelmezett) Binance-adaton szimulált végrehajtással dolgozik. A
visszateszt-mód ugyanazt a pipeline-t futtatja historikus adatokon, Binance
kapcsolat nélkül, \texttt{adapters/testing/backtest.py} segítségével;
eredményeit a konzolra és CSV-be írja. Az offline tanítás a teljes dataset-et
(alapértelmezetten 3 nap BTCUSDT 5 perces gyertyái) egy menetben dolgozza fel,
majd a súlyokat \texttt{*.pt} fájlba menti. A \emph{csak-dashboard} mód végül
akkor hasznos, ha egy korábbi futás eredményeit akarjuk áttekinteni: ekkor
a Binance-feed kikapcsolható, és a dashboard offline módban is böngészhető.

\subsubsection*{Futásidei naplózás és leállítás}

A \texttt{stdout} a rendes állapotüzeneteket tartalmazza
(\texttt{[BOOT]}, \texttt{[SIGNAL]}, \texttt{[FILL]}, \texttt{[DASHBOARD]}),
a \texttt{stderr} pedig a figyelmeztetéseket és hibákat
(\texttt{[WARN]}, \texttt{[FATAL ERROR]}). A két csatorna külön is
naplózható:
\begin{verbatim}
nohup python main.py > run_log.txt 2> run_err.txt &
\end{verbatim}

A futtatás \textbf{Ctrl+C}-vel bármikor biztonságosan megszakítható.
A \texttt{finally} ág garantálja, hogy (i) a WebSocket-kapcsolat lezárul,
(ii) az in-memory esemény-pufferek kiürülnek,
(iii) a HTTP-szerver lezárja az aktív kapcsolatokat, és
(iv) a végleges állapot (nyitott pozíciók, tőkeállás, utolsó jelzés)
kiírásra kerül a \texttt{FINAL DASHBOARD SNAPSHOT} fejléc alatt.

\subsubsection*{Biztonsági megfontolások}

A demó alapértelmezés szerint \emph{paper-trading} módban fut, valódi
Binance-rendelést nem küld. Ha mégis éles kulcs kerül a \texttt{.env} fájlba,
a Binance-oldalon ajánlott (a) csak olvasási és Futures-kereskedési
jogosultságot adni, és (b) IP-whitelisttel megszorítani a kulcs használatát.
Az irányítópult alapértelmezetten a \texttt{0.0.0.0:8080} címen hallgat;
nyilvános IP-n csak külső \texttt{reverse proxy} (nginx, traefik) mögött
ajánlott közzétenni.

\subsection{Hibaelhárítás}

\noindent\textbf{\texttt{ModuleNotFoundError: mamba\_ssm}}\\
Telepítsd a \texttt{requirements-ml.txt}-ben szereplő CUDA-kompatibilis csomagokat,
vagy futtasd CPU-n: \texttt{python main.py --device cpu}.

\noindent\textbf{Binance API 403/401}\\
Ellenőrizd a \texttt{.env} fájl kulcsait és az
IP-whitelist beállításait a Binance kereskedési felületén.

\noindent\textbf{WebSocket kapcsolódási hiba}\\
A tűzfal blokkolhatja a Binance \texttt{wss://}
végpontját ($9443$ port); szükség esetén proxyt kell beállítani.

% ─────────────────────────────────────────────────────────────────────────────
\section{Tesztelés}
\label{sec:testing}
% ─────────────────────────────────────────────────────────────────────────────

\subsection{Tesztelési stratégia}

A tesztkészlet a \textbf{pytest} keretrendszert használja. A core logikát
egységtesztek fedik le; a Binance és aiohttp adapterek \texttt{unittest.mock}-kal
vannak kiváltva, tehát hálózat nélkül, determinisztikusan futnak.

\begin{table}[H]
  \centering
  \small
  \begin{tabular}{|l|l|p{7cm}|}
    \hline
    \textbf{Tesztfájl} & \textbf{Típus} & \textbf{Mit fedez le?} \\
    \hline\hline
    \texttt{test\_strategy.py}          & Egységteszt   & ThresholdStrategy jelzéslogika, küszöb-, konfigurációs- és méretezetségi ellenőrzések \\
    \texttt{test\_risk\_policy.py}       & Egységteszt   & RiskPolicy blokkolási útvonalak, MaxQty, MaxNotional, Confidence, SlippageGuard szabályok \\
    \texttt{test\_execution.py}         & Integrációs   & PaperBroker megbízás-végrehajtás, kitöltési szimuláció, Stop és Limit megbízások \\
    \texttt{test\_equity.py}            & Egységteszt   & SimulationWallet tőkeszámítás, PnL, mark-price frissítés \\
    \texttt{test\_fee\_fix.py}          & Regressziós   & Jutalékszámítás pontossága kereskedési cikluson át \\
    \texttt{test\_fill\_fix.py}         & Regressziós   & Fill-esemény idő- és árlogikája, részleges kitöltés kezelése \\
    \texttt{test\_signal\_diagnostics.py} & Egységteszt & Jelzés-diagnosztikai metrikák (edge score, hit-rate számítás) \\
    \texttt{test\_binance\_feed.py}     & Mock-alapú    & BinanceFeed WebSocket-esemény feldolgozás, reconnect logika \\
    \hline
  \end{tabular}
  \caption{A rendszer tesztfájljainak összefoglalása.}
  \label{tab:tests}
\end{table}

\subsection{Tesztek futtatása}

\begin{verbatim}
# Összes teszt futtatása
pytest tests/ -v

# Csak az egységtesztek, kihagyva a lassú integrációs teszteket
pytest tests/ -v -m "not slow"

# Lefedettségi riport generálása
pytest tests/ --cov=core --cov-report=term-missing
\end{verbatim}

\subsection{Tesztesetek és eredmények}

Az alábbiakban néhány jellemző teszteset kerül bemutatásra.

\begin{itemize}
    \item A \texttt{test\_risk\_policy.py} ellenőrzi, hogy az összes
          tiltó szabály (\texttt{MaxQtyRule}, \texttt{MaxNotionalRule},
          \texttt{ConfidenceRule}, \texttt{SlippageGuardRule}) helyesen
          blokkol érvénytelen megbízásokat, és hogy a
          \texttt{RiskPolicy} együttes alkalmazásakor a legszigorúbb
          szabály érvényesül. A tesztek véletlenszerű bemenetek mellett
          ($\sim 100$ kombináció) is megbízhatóan lefedik a határeseteket.
    \item A \texttt{test\_strategy.py} verifikálja, hogy a
          \texttt{ThresholdStrategy} pontosan akkor ad
          \textsc{Buy}/\textsc{Sell} jelzést, ha az előrejelzett
          valószínűség meghaladja a beállított küszöböt
          (\texttt{threshold}), és hogy a \texttt{min\_edge},
          \texttt{min\_confidence} és \texttt{max\_sigma} paraméterek
          megfelelően szűrik a gyenge jelzéseket.
    \item A \texttt{test\_equity.py} és \texttt{test\_fill\_fix.py}
          összes tesztje zöld, ami megerősíti, hogy a PnL-számítás és a
          részleges kitöltések kezelése helyesen működik
          multi-szimbólumos portfóliókban.
    \item A \texttt{test\_fee\_fix.py} regressziós tesztként
          garantálja, hogy egy korábbi bugfix (jutalék duplán számolása
          tételes nézeten) nem kerül vissza a kódbázisba.
\end{itemize}

\subsection{Tesztelési korlátok és fejlesztési lehetőségek}

Nincs end-to-end integrációs teszt, amely a teljes pipeline-t egyetlen
futtatásban ellenőrizné. A visszatesztelési mód részben betölti ezt a szerepet,
de egy dedikált szimulációs tesztsuite még jobb fedettséget adna.
